% Original PDF 22 / printed 12, complete Section 2.2.
% Fraktur source Z/B use the explicitly user-approved non-Fraktur mathcal rendering.
% block: cycles-boundaries-heading page:22 type:section
\subsection{{Cycles, boundaries, and homology}}\label{sec:cycles-boundaries-homology}

% block: cycles-boundaries-homology-definition page:22 type:ordinary-paragraph
{Given a chain complex $(V,\partial)$, the null space $\mathcal Z$ and the range $\mathcal B$ of $\partial$ are each graded subspaces of $V$. (They are subcomplexes, but with pretty boring differentials, since the restriction of $\partial$ to each is zero.) The elements of $\mathcal Z_k$ are called the $k$-cycles and the elements of the range $\mathcal B_k$ of $\partial_{k+1}$ are called the $k$-boundaries. The differential property $\partial^2=0$ means that $\mathcal B_k\subset\mathcal Z_k$, and the $k$th \emph{homology space} is defined to be the quotient space $\mathcal H_k=\mathcal Z_k/\mathcal B_k$. Thus the elements of the homology space are equivalence classes of cycles, with two $k$-cycles considered equivalent if their difference is the boundary of a $(k+1)$-chain. If the boundary spaces $\mathcal B_k$ are exactly equal to the cycle spaces $\mathcal Z_k$, rather than proper subspaces, then the homology spaces $\mathcal H_k$ all vanish. In this case the complex is said to be \emph{exact}.}
